\section{Data records}\label{records}

The source and derived-data package contains graph modules, curated
tables, source and correction ledgers, validation evidence and workflow
code. The machine-readable manifest
\path{PRE_RELEASE_MANIFEST.tsv} identifies payload sizes, checksums,
categories and release/licence status. Separately retrieved bulk
payloads are listed in the bulk manifest. Data Availability reports
the current access routes and outstanding sequence-release gate.

KG v3.0.0 contains 45,706,977 asserted triples from 16 manifested
modules. The website distributes three versioned artifacts: the
scientific source package, the original RDF modules and the
deduplicated asserted graph in compressed N-Quads format.
The \href{https://rubalkhali.science/downloads/kg/3.0.0/manifest.json}{release manifest}
identifies their contents, sizes and checksums.

\subsection{Knowledge-graph modules}

The primary artifacts are the knowledge-graph modules, distributed as
RDF/XML and Turtle files and structured by data type
(Table~\ref{tab:modules}). The base ontology
(\texttt{rubalkhali.owl}) holds the project terminology. Domain modules
share site, specimen and process identifiers; the release also includes
the generated taxonomic vocabulary and four pinned reference ontologies.
These modules load together into an RDF store. Domain queries require the corresponding
shared terminology and linked site, specimen or process modules.

\begin{table}[h]
\centering
\caption{The 16 RDF modules included in KG v3.0.0.}
\label{tab:modules}
\small
\begin{tabular}{p{0.38\textwidth} p{0.54\textwidth}}
\toprule
\textbf{Module} & \textbf{Content} \\
\midrule
\texttt{rubalkhali.owl} & Base ontology (TBox) \\
\texttt{rubalkhali\_sites.owl} & Sampling sites with coordinates and biome annotations \\
\texttt{rubalkhali\_samples.owl} & Soil and plant specimen individuals with compartment types \\
\texttt{rubalkhali\_measurements.owl} & Field environmental measurements and climate records \\
\texttt{rubalkhali\_xrf.owl} & XRF measurement values from the field and laboratory workflows \\
\texttt{rubalkhali\_dna.owl} & DNA extracts with concentration data \\
\texttt{rubalkhali\_sra.owl} & Sequencing run records linked to ENA accessions \\
\texttt{rubalkhali\_qc.owl} & Sequencing quality metrics (GC content, duplicate rate, sequence counts) \\
\texttt{rubalkhali\_controls.ttl} & Laboratory-control materials, roles, batches and sequence occurrences \\
\path{rubalkhali_ph_eq_ph_shared_v1_0_1.ttl} & Archived-soil pH measurements, sessions and reconciled specimen links \\
\texttt{rubalkhali\_taxonomy\_abox.ttl} & Taxon observations over feature profiles \\
\texttt{ecosystem\_module.ttl} & Vocabulary supporting canonical taxonomic mappings \\
\texttt{sio.owl} & Semanticscience Integrated Ontology v1.59 \\
\texttt{envo.owl} & Environment Ontology, 20 October 2025 \\
\texttt{pato.owl} & Phenotype and Trait Ontology, 14 May 2025 \\
\texttt{uo.owl} & Units Ontology, 25 May 2023 \\
\bottomrule
\end{tabular}
\end{table}

\subsection{Tabular data products}

Alongside the graph, the deposit distributes the curated tabular data
products from which the modules are generated: the versioned source
ledger linking site and trip to specimen, extraction, library,
sequence run, quality-control decision and derived community; the
canonical Trips~1--5 amplicon feature table and its taxonomy mapping;
the correction and alias ledgers that record every deterministic
metadata repair without overwriting source values (the alias ledger
maps the 36 genuine Trip-1-only records of numeric sites 61--64 to the
four existing named site individuals); the curated field
environmental table; the field and laboratory XRF tables together with
the audited chemical-mapping ledger; and the versioned archived-soil
pH dataset \texttt{EQ-PH-SHARED-v1.0.1}. The canonical feature table
holds 1,271 profiles, of which the ecological analysis retains 1,237
(1,227 profiles from the primary frame of sites~1--60); the 34 excluded
profiles are the 24 controls, nine biological profiles
below the declared 1,000-read ecological threshold and one additional
\texttt{T1Dr1} run with 934 reads. The four derived
companion-analysis inputs are distributed as records for downstream
statistical re-execution: the paired PMA
table (\path{metadata/relic-dna/PMA_ASV_table.tsv}; 18 analytical
columns from nine aliquot pairs, PMA-treated and untreated, with
controls kept separately), the 150 CoverM profiles
(\path{metadata/metagenome/coverm_profiles.tar.gz}), the eggNOG
annotations (\path{metadata/metagenome/eq.emapper.annotations.gz})
behind the 990-genome analytical subset, and the measured-function
input archive (\path{metadata/metagenome/measured_function_inputs.tar.gz})
with the six exact source tables of the genome-derived versus PICRUSt2
comparison. These are fixed inputs for downstream analyses. Methods identifies
the additional raw reads, upstream workflows and acquisition metadata
needed for end-to-end reproduction.

\subsection{Validation schemas, queries and evidence}

The deposit includes the artifacts needed to re-validate the resource:
the ShEx shapes that define the structural contract for each entity
family, the positive and negative structural fixtures, the archived
competency SPARQL queries, and the machine-readable evidence bundles
(reconciliation summaries, audit tables, validation records and
checksum manifests) produced by the release workflow. A companion data
dictionary (\path{metadata/DATA_DICTIONARY.tsv}) records the field
names, data types, units, missing-value conventions and definitions
for the curated canonical tabular records. Quantities carry units only
where a unit is documented at the source: environmental source tables
retain degrees Celsius, hectopascals (millibars) and percent relative
humidity; in the graph, pressure magnitudes are
multiplied by 100 and asserted in pascals (\texttt{obo:UO\_0000110}),
climate
records use degrees Celsius and millimetres, and DNA concentrations
are in nanograms per microlitre. XRF values are instrument-reported magnitudes with undocumented
concentration units.

The manifests and data dictionary provide the versioned file-level
inventory. The final archival dataset citation remains a release gate
(Section~\ref{data_availability}).
